\documentclass[11pt,a4paper]{ltjsarticle}
\usepackage{amsmath,amssymb,amsthm,mathtools}
\usepackage{hyperref}
\newtheorem{theorem}{定理}[section]
\newtheorem{proposition}[theorem]{命題}
\newtheorem{remark}[theorem]{注意}
\renewcommand{\proofname}{証明}
\newcommand{\C}{\mathbb C}
\newcommand{\PP}{\mathbb P}
\newcommand{\Spec}{\operatorname{Spec}}
\newcommand{\Sym}{\operatorname{Sym}}
\newcommand{\Gr}{\operatorname{Gr}}
\title{三次三様体の Poisson 錐と二次ジェット上の量子化}
\author{}
\date{2026年11月10日・研究結果のソース整理版}
\begin{document}
\maketitle
\begin{abstract}
滑らかな三次三様体の Poisson 錐が Pl\"ucker 埋め込みされた
$\Gr(2,5)$ のアフィン錐であることを二次方程式のイデアルとして述べ，
その頂点の二次無限小近傍上の平坦な次数付き量子化を
二次 Kontsevich 積により明示する．
本稿は2026年11月3日の Poisson ジェット原稿で扱った主張のうち，
独立した二次ジェットの議論を自立したソースとして保持する．
有限係数量子化，17次元の一般族および ambient の復元については
\texttt{papers/generic17\_ambient\_recovery.tex} と
\texttt{papers/sections/poisson\_geometry.tex} に証明を置く．
\end{abstract}
\section{Poisson 錐の二次方程式}
$W$ を五次元複素ベクトル空間，
$S=\Sym W=\C[x_0,\ldots,x_4]$ とし，
$F\in S_3$ が滑らかな三次三様体
$X=V(F)\subset\PP(W^*)$ を定めるとする．
随伴公式と三次元における自然同型によって
\[
K_X\simeq\mathcal O_X(-2),\qquad
\bigwedge^2T_X\simeq\Omega_X^1(2),\qquad
H^0(X,\bigwedge^2T_X)\simeq\bigwedge^2W
\]
となる．最後の同型は Euler 列と conormal 列から従う．
$p=(p_{ij})\in\bigwedge^2W$ に対して
\[
\beta_p=\sum_{i<j}p_{ij}\,dx_i\wedge dx_j,\quad
E=\sum_{i=0}^4x_i\partial_{x_i},\quad
\alpha_p=\iota_E\beta_p,\quad
\Omega=dx_0\wedge\cdots\wedge dx_4
\]
とおく．定数ベクトル $v_p$ を
$\beta_p\wedge\beta_p=\iota_{v_p}\Omega$
で定義する．その成分 $q_i(p)$ は五つの Pl\"ucker 二次式である．

\begin{theorem}[二次方程式のスキームとしての Poisson 錐]
\label{thm:cone}
Poisson 条件を Schouten 括弧で与えたアフィンスキームは
\[
\operatorname{Pois}(X)
=\Spec\frac{\C[p_{ij}\mid0\le i<j\le4]}
{(q_0,\ldots,q_4)}
=\widehat{\Gr}(2,W)
\]
である．とくに非零部分の射影化は $\Gr(2,5)$，
錐は既約かつ七次元である．
\end{theorem}
\begin{proof}
$d\alpha_p=2\beta_p$ より
$\alpha_p\wedge d\alpha_p=\iota_E(\beta_p\wedge\beta_p)$ であり，
Euler 恒等式 $E(F)=3F$ から
\[
\alpha_p\wedge d\alpha_p\wedge dF
=v_p(F)\,\iota_E\Omega-3F\,\iota_{v_p}\Omega
\tag{1}
\]
を得る（体積形式の縮約の符号は上式で固定する）．
$X$ 上の Frobenius 条件は $v_p(F)$ が $S/(F)$ で零になること
と同値である．これは二次式なので $S_2$ で零になることと同値である．
線形写像 $W^*\to S_2,\ v\mapsto v(F)$ は単射である．
実際 $v(F)=0$ で $v\ne0$ なら $F$ は座標変換後に一変数に依存せず，
$X$ は特異な射影錐になる．
従って (1) の係数比較から五つの $q_i(p)$ 自身が
Poisson 方程式の生成元である．単なる集合論的同一視ではない．
\end{proof}

\section{頂点の二次無限小近傍}
$\mathfrak m=(p_{ij})$ とおいて
\[
R_2=
\frac{\C[p_{ij}]}
{(q_0,\ldots,q_4)+\mathfrak m^3}
\]
を考える．これは Poisson 錐の頂点の二次無限小近傍の座標環である．
同じ $F$ に対し五変数アフィン空間上の次数二の双ベクトル場
$\Pi_p$ を
\[
\iota_{\Pi_p}\Omega=dF\wedge\beta_p
\tag{2}
\]
で定める．係数は $p_{ij}$ に線形であり，$F$ は
$\Pi_p$ の Casimir である．$\Pi_p$ は Pl\"ucker 方程式の下で
Jacobi 恒等式を満たす．
添字の総和規約で $f_i=\partial_i f$ 等と書き，$S\otimes R_2$ 上の積を
\[
\begin{aligned}
f\star g={}&fg+\frac12\Pi^{ij}f_i g_j
+\frac18\Pi^{ij}\Pi^{k\ell}f_{ik}g_{j\ell}\\
&+\frac1{12}\Pi^{ij}(\partial_j\Pi^{k\ell})
  (f_{ik}g_\ell-f_k g_{i\ell})
\end{aligned}
\tag{3}
\]
で定める．これは Kontsevich の二次の星積を
$\mathfrak m^3=0$ で打ち切ったものである．
\begin{theorem}[普遍二次ジェットの平坦量子化]
$\star$ は結合的な次数付き $R_2$-代数構造を定め，
$F$ はその中心に属する．商
\[
\mathcal A_2=(S\otimes R_2,\star)/(F)
\]
は次数ごとに $R_2$-自由であり，
特別ファイバーは $S/(F)$ である．その一次の交換子は
$\Pi_p$ の Poisson 括弧に一致する．
\end{theorem}
\begin{proof}
(3) の結合性を二次まで計算した際の障害は
Jacobi 恒等式 $[\Pi_p,\Pi_p]=0$ の普遍的線形結合である．
定理\ref{thm:cone}によりこれは $R_2$ 上で零で，
$\mathfrak m^3=0$ より高次項も零である．
一次の交換子は $\Pi^{ij}f_i g_j$ である．
二次項は対称なので交換子の二次部分は消える．
特に (2) により $\Pi_p(dF,-)=0$ だから，
$F$ は $\star$-中心的である．
各次数で $F\star-$ は通常の単射
$S_{d-3}\to S_d$ の冪零摂動である．
その余核は $R_2$ 上の自由加群に同型である．
実際，多項式環の各次数でこの単射の像に
$\C$-線形補空間を選べば，
$\mathfrak m$ に関する冪零行列の消去により
$\operatorname{im}(F\star-)$ の補空間も同じ rank で選べる．
従って $\mathcal A_{2,d}$ は $R_2$-自由である．
\end{proof}

\begin{remark}[意味と限界]
頂点そのもののファイバーは可換な $X$ である．
非被約な二次近傍上の族は，頂点から出る一次方向に対して
$p\wedge p=0$ が二次障害になることを記録する．
この主張は，七次元の被約錐全体上で有限係数の普遍族が存在する
という主張ではない．一般の非零方向に沿う有限係数量子化と
中心三次超曲面および17次元の真の変形については，
リポジトリの別原稿で扱う．
\end{remark}

\section{既知の結果との関係と再生成}
本稿で用いる形式量子化の普遍的な二次積は
Kontsevich の既知の公式である．固定した滑らかな三次三様体の
非零 Poisson 構造の集合論的分類は
Loray--Pereira--Touzet と Mayanskiy に依存する．
本稿の二次方程式の計算はスキーム論的な等式を明記したものである．
既存の有限パラメータでの Sklyanin 型代数，
全混合項の中心三次元，および ambient 非可換 $\PP^4$ の復元は
\texttt{papers/generic17\_ambient\_recovery.tex} と
\texttt{papers/sections/poisson\_geometry.tex} を参照する．
本稿はこれらの長い証明を重複しない．
原稿の再生成には LuaLaTeX と LuaTeX-ja を用いる．

\begin{thebibliography}{9}
\bibitem{Kontsevich}
M.~Kontsevich,
Deformation quantization of Poisson manifolds,
Lett. Math. Phys. \textbf{66} (2003), 157--216.
\bibitem{LPT}
F.~Loray, J.~V.~Pereira and F.~Touzet,
Foliations with trivial canonical bundle on Fano 3-folds,
Math. Nachr. \textbf{286} (2013), 921--940.
\bibitem{Mayanskiy}
E.~Mayanskiy,
Poisson cohomology of two Fano threefolds,
\url{https://arxiv.org/abs/1303.5916}.
\end{thebibliography}
\end{document}
